\documentclass[11pt,a4paper]{article}  % Déclare la classe du document.

\newcommand{\chaptermark}[2]{\og #2 \fg{} (#1)} % ligne uniquement pour éviter une erreur en cas d'article
\usepackage[table,dvipsnames,svgnames]{xcolor}
\usepackage{listingsutf8}
%\usepackage[utf8]{inputenc}% gestion des accents (source)
\usepackage[T1]{fontenc}% gestion des accents (PDF)
\usepackage{fontspec}
\usepackage[french]{babel}% gestion du français
\usepackage{textcomp}% caractères additionnels
\usepackage{mathtools,amssymb,amsthm}% packages de l'AMS + mathtools
\usepackage{amsmath}
\usepackage{yhmath}
\usepackage{mathpazo}
\usepackage[scaled=0.95]{helvet}
\usepackage{courier}
%\usepackage{fourier}
%\usepackage[scaled=0.95]{helvet}
%\usepackage{courier}
\usepackage{bm}
\usepackage{gensymb}
\usepackage{graphicx}

\usepackage{enumitem}
\usepackage{pifont}

\usepackage{pstricks}
\usepackage{pstricks-add}
\usepackage{pst-eucl}
\usepackage{pgfplots}
\let\clipbox\relax
%\usepackage{pst-plot,pst-text,pst-tree,pstricks-add}
%\usepackage{pst-node} % required package
\usepackage{pst-all}
\usepackage[most,skins,breakable]{tcolorbox}

\usepackage{euscript}





% ############### Nouvelles commandes ###############
% Coloration des labels des points dans pst-eucl
\makeatletter
\define@key[psset]{pst-eucl}{LabelColor}{\def\psk@LabelColor{#1}}
\psset{LabelColor=black}
\def\Pst@WhichLabel#1{{\color{\psk@LabelColor}%
		\ifx\psk@PointName\@default#1\else\psk@PointName\fi}}
\makeatother

% renomme le symbole #2 en #1#2 et invalide le symbole initial #2:
\def\renamesymbol#1#2{
	\expandafter\let\expandafter\newsym\expandafter=\csname#2\endcsname
	\expandafter\global\expandafter\let\csname#1#2\endcsname=\newsym
	\expandafter\global\expandafter\let\csname#2\endcsname=\origsym
}



\usepackage{ProfCollege}
\renamesymbol{Casio}{Calculatrice}
\usepackage[nonshellescape]{ProfLycee}



%\usepackage[left=1.5cm,right=1.5cm,top=2.5cm,bottom=2.5cm]{geometry}
\usepackage[left=13mm,right=13mm,top=18mm,bottom=18mm,footskip=30pt]{geometry}
\setlength{\headheight}{15pt}

\usepackage{titlesec}
\usepackage{pas-math}
\usepackage{longtable}
\usepackage{tikz,tkz-tab}
\usepackage{amsmath}
\usepackage[right]{eurosym}
\usepackage{sectsty}
\usepackage{tipfr}
\usepackage{setspace}

%pour faire les touches de la calculatrices
\usepackage{menukeys}
\renewmenumacro{\menu}{roundedmenus} % default: menus
\renewmenumacro{\directory}{pathswithfolder} % default: paths
\renewmenumacro{\keys}{shadowedroundedkeys} % default: roundedkeys


\usepackage{multicol}
\usepackage{multirow}
\usepackage{systeme}

%% Format des titres de chapitres, paragraphes, ...
\titleformat{\chapter}[frame]
{\color{red!85!white}\normalsize}%
{\filright\sffamily\Large%
	\enspace Chapitre \thechapter\enspace}%
{8pt}
{\color{red!85!white}\rule{0pt}{30pt}\sffamily\Huge\bfseries\filcenter}
%%%%%%%%%%%%%%%%%%%%
\titlespacing{\chapter}{0pc}{0pc}{2pc}
\titlespacing{\subsection}{2pc}{}{}

\titleformat{\subsection}[block]{\hspace{1em}}{\thesubsection}{1em .}{}
\titleformat{\subsubsection}[block]{\hspace{2em}}{\thesubsection}{1em )}{}

\chapterfont{\color{red}{}}
\sectionfont{\color{green!55!black}{}}
\subsectionfont{\color{cyan!65!black}{}}
\subsubsectionfont{\color{violet}{}}
\renewcommand{\thesection}{\color{green!55!black}\Roman{section}.}
\renewcommand{\thesubsection}{\color{cyan!65!black}\arabic{subsection})}
\renewcommand{\thesubsubsection}{\color{violet}\alph{subsubsection})}
\setcounter{secnumdepth}{3}

\titleformat{\subsection}[block]{\hspace{1em}}{\thesubsection}{1em .}{}
\titleformat{\subsubsection}[block]{\hspace{2em}}{\thesubsection}{1em )}{}

\usepackage[column=A]{cellspace}
\setlength{\cellspacebottomlimit}{4pt}
\setlength{\cellspacetoplimit}{4pt}
\setlength{\parindent}{2pt}
\setlength{\tabcolsep}{0.5cm}


%entête et pieds de pages
\usepackage{fancyhdr}
\pagestyle{fancy}
\fancyhf{}
\usepackage{lastpage}
\renewcommand{\chaptermark}[1]{\markboth {Chap. \thechapter~ - ~#1}{}}
\renewcommand\headrulewidth{1pt}
\fancyhead[LO,RE]{1ère G - spécialité maths}
\fancyhead[RO,LE]{\leftmark}
\renewcommand\footrulewidth{1pt}
\fancyfoot[LO,RE]{Année 2025-2026}
\fancyfoot[RO,LE]{\textbf{Page \thepage/\pageref{LastPage}}}

%Mes boites pour les définitions, propriétés, remarques, ...

%Mes boites pour les définitions, propriétés, remarques, ...

\newcommand\titrebox[1]{\textbf{#1}}
\newtcolorbox{definition}{breakable,colback=red!10!white,
	colframe=red!85!white, fonttitle=\bfseries,enhanced,
	attach boxed title to top left={xshift=5mm, yshift=-2mm},boxed title style={size=small,colback=red!10!white,colframe=red!85!white, arc=0.5mm},
	title=\textcolor{red!85!white}{Définition} , arc=3mm}

\newtcolorbox{definitionTitre}[1]{breakable,colback=red!10!white,
	colframe=red!85!white, fonttitle=\bfseries,enhanced,
	attach boxed title to top left={xshift=5mm, yshift=-2mm},boxed title style={size=small,colback=red!10!white,colframe=red!85!white, arc=0.5mm},
	title=\textcolor{red!85!white}{Définition - #1} , arc=3mm}

\newtcolorbox{definitionPerso}[1]{breakable,colback=red!10!white,
	colframe=red!85!white, fonttitle=\bfseries,enhanced,
	attach boxed title to top left={xshift=5mm, yshift=-2mm},boxed title style={size=small,colback=red!10!white,colframe=red!85!white, arc=0.5mm},
	title=\textcolor{red!85!white}{#1} , arc=3mm}

\newtcolorbox{regle}{breakable,colback=red!10!white,
	colframe=red!85!white, fonttitle=\bfseries,enhanced,
	attach boxed title to top left={xshift=5mm, yshift=-2mm},boxed title style={size=small,colback=red!10!white,colframe=red!85!white, arc=0.5mm},
	title=\textcolor{red!85!white}{Règle} , arc=3mm}

\newtcolorbox{regles}{breakable,colback=red!10!white,
	colframe=red!85!white, fonttitle=\bfseries,enhanced,
	attach boxed title to top left={xshift=5mm, yshift=-2mm},boxed title style={size=small,colback=red!10!white,colframe=red!85!white, arc=0.5mm},
	title=\textcolor{red!85!white}{Règles} , arc=3mm}

\newtcolorbox{definitions}{breakable,colback=red!10!white,
	colframe=red!85!white, fonttitle=\bfseries,enhanced,
	attach boxed title to top left={xshift=5mm, yshift=-2mm},boxed title style={size=small,colback=red!10!white,colframe=red!85!white, arc=0.5mm},
	title=\textcolor{red!85!white}{Définitions} , arc=3mm}

\newtcolorbox{propriete}{breakable,colback=red!10!white,
	colframe=red!85!white, fonttitle=\bfseries,enhanced,
	attach boxed title to top left={xshift=5mm, yshift=-2mm},boxed title style={size=small,colback=red!10!white,colframe=red!85!white, arc=0.5mm},
	title=\textcolor{red!85!white}{Propriété} , arc=3mm}

\newtcolorbox{proprieteTitre}[1]{breakable,colback=red!10!white,
	colframe=red!85!white, fonttitle=\bfseries,enhanced,
	attach boxed title to top left={xshift=5mm, yshift=-2mm},boxed title style={size=small,colback=red!10!white,colframe=red!85!white, arc=0.5mm},
	title=\textcolor{red!85!white}{Propriété - #1} , arc=3mm}

\newtcolorbox{proprietes}{breakable,colback=red!10!white,
	colframe=red!85!white, fonttitle=\bfseries,enhanced,
	attach boxed title to top left={xshift=5mm, yshift=-2mm},boxed title style={size=small,colback=red!10!white,colframe=red!85!white, arc=0.5mm},
	title=\textcolor{red!85!white}{Propriétés} , arc=3mm}

\newtcolorbox{theoreme}{breakable,colback=red!10!white,
	colframe=red!85!white, fonttitle=\bfseries,enhanced,
	attach boxed title to top left={xshift=5mm, yshift=-2mm},boxed title style={size=small,colback=red!10!white,colframe=red!85!white, arc=0.5mm},
	title=\textcolor{red!85!white}{Théorème} , arc=3mm}

\newtcolorbox{theoremeTitre}[1]{breakable,colback=red!10!white,
	colframe=red!85!white, fonttitle=\bfseries,enhanced,
	attach boxed title to top left={xshift=5mm, yshift=-2mm},boxed title style={size=small,colback=red!10!white,colframe=red!85!white, arc=0.5mm},
	title=\textcolor{red!85!white}{Théorème - #1} , arc=3mm}

\newtcolorbox{theoremes}{breakable,colback=red!10!white,
	colframe=red!85!white, fonttitle=\bfseries,enhanced,
	attach boxed title to top left={xshift=5mm, yshift=-2mm},boxed title style={size=small,colback=red!10!white,colframe=red!85!white, arc=0.5mm},
	title=\textcolor{red!85!white}{Théorèmes} , arc=3mm}

\newtcolorbox{bexemple}{
	enhanced,
	breakable,
	fonttitle=\bfseries,
	frame empty,
	nobeforeafter,
	borderline west={0.5mm}{0mm}{cyan!75!black},
	title=\textcolor{cyan!75!black}{\faHubspot\:\:Exemple},
	%title after break=Exemple (suite),
	colback=white,
	extras={enhanced, frame empty},
	detach title,before upper={\tcbtitle\ :\quad },
}

\newtcolorbox{bexemples}{
	enhanced,
	breakable,
	fonttitle=\bfseries,
	frame empty,
	nobeforeafter,
	borderline west={0.5mm}{0mm}{cyan!75!black},
	title=\textcolor{cyan!75!black}{\faHubspot\:\:Exemples},
	%title after break=Exemples (suite),
	colback=white,
	extras={enhanced, frame empty},
	detach title,before upper={\tcbtitle\ :\quad },
}

\newtcolorbox{bexemplet}[1]{
	enhanced,
	breakable,
	fonttitle=\bfseries,
	frame empty,
	borderline west={0.5mm}{0mm}{cyan!75!black},
	title=\textcolor{cyan!75!black}{\faHubspot\:\:Exemple - #1},
	colback=white,
	extras={enhanced, frame empty},
	detach title,before upper={\tcbtitle\ :\quad },
}

\newtcolorbox{brmq}{
	enhanced,
	breakable,
	fonttitle=\bfseries,
	frame empty,
	nobeforeafter,
	borderline west={0.5mm}{0mm}{orange!95!white},
	title=\textcolor{orange!85!white}{\faHandPointRight[regular]\:\:Remarque} ,
	colback=white,
	extras={enhanced, frame empty},
	detach title,before upper={\tcbtitle\ \quad },
}

\newtcolorbox{brmqs}{
	enhanced,
	breakable,
	fonttitle=\bfseries,
	frame empty,
	nobeforeafter,
	borderline west={0.5mm}{0mm}{orange!85!white},
	title=\textcolor{orange!85!white}{\faHandPointRight[regular]\:\:Remarques},
	colback=white,
	detach title,before upper={\tcbtitle\ \quad },
	extras={enhanced, frame empty},
}

\newtcolorbox{brmqt}[1]{
	enhanced,
	breakable,
	fonttitle=\bfseries,
	frame empty,
	nobeforeafter,
	borderline west={0.5mm}{0mm}{orange!85!white},
	title=\textcolor{orange!85!white}{\faHandPointRight[regular]\:\:Remarque - #1},
	colback=white,
	extras={enhanced, frame empty},
	detach title,before upper={\tcbtitle\ \quad },
}

\newtcolorbox{consequences}[1]{
	enhanced,
	breakable,
	fonttitle=\bfseries,
	frame empty,
	nobeforeafter,
	borderline west={0.5mm}{0mm}{ProcessBlue},
	title=\textcolor{ProcessBlue}{#1},
	colback=white,
	extras={enhanced, frame empty},
	detach title,before upper={\tcbtitle\ :\quad },
}

\newtcolorbox{bmethode}{
	enhanced,
	breakable,
	fonttitle=\bfseries,
	frame empty,
	nobeforeafter,
	borderline west={0.5mm}{0mm}{blue!75!black},
	title=\textcolor{blue!75!black}{Méthode} ,
	colback=white,
	extras={enhanced, frame empty},
	attach boxed title to top left = {xshift = -3mm},
	boxed title style={size=small,colback=white,colframe=blue!75!black, arc=3mm,
	}
}

\newtcolorbox{bmethodes}{
	enhanced,
	breakable,
	fonttitle=\bfseries,
	frame empty,
	nobeforeafter,
	borderline west={0.5mm}{0mm}{blue!75!black},
	title=\textcolor{blue!75!black}{Méthodes} ,
	colback=white,
	extras={enhanced, frame empty},
	attach boxed title to top left = {xshift = -3mm},
	boxed title style={size=small,colback=blue!75!black,colframe=blue!75!black, arc=3mm,
	}
	boxed title after break style={size=small,colback=white,colframe=blue!75!black, arc=3mm,
	}
}

\newtcolorbox{bmethodet}[1]{
	enhanced,
	breakable,
	fonttitle=\bfseries,
	frame empty,
	nobeforeafter,
	borderline west={0.5mm}{0mm}{blue!75!black},
	title=\textcolor{blue!75!black}{\faTasks\:\:Méthode - #1} ,
	colback=white,
	extras={enhanced, frame empty},
	attach boxed title to top left = {xshift = -3mm},
	boxed title style={size=small,colback=white,colframe=blue!75!black, arc=3mm,
	}
}

\newtcolorbox{bdemo}{
	enhanced,
	breakable,
	fonttitle=\bfseries,
	frame empty,
	nobeforeafter,
	borderline west={0.5mm}{0mm}{violet!65!white},
	title={\faEdit[regular]\:\:Démonstration},
	%title after break=Démonstration (suite),
	colback=white,
	extras={enhanced, frame empty},
	attach boxed title to top left = {xshift = -3mm},
	boxed title style={size=small,colback=violet!65!white,colframe=violet!65!white, arc=3mm,
	}
}

\newtcolorbox{bdemos}{
	enhanced,
	breakable,
	fonttitle=\bfseries,
	frame empty,
	nobeforeafter,
	borderline west={0.5mm}{0mm}{violet!65!white},
	title={\faEdit[regular]\:\:Démonstrations},
	%title after break=Démonstration (suite),
	colback=white,
	extras={enhanced, frame empty},
	attach boxed title to top left = {xshift = -3mm},
	boxed title style={size=small,colback=violet!65!white,colframe=violet!65!white, arc=3mm,
	}
}

\newtcolorbox{bdemot}[1]{
	enhanced,
	breakable,
	fonttitle=\bfseries,
	frame empty,
	nobeforeafter,
	borderline west={0.5mm}{0mm}{violet!65!white},
	title={\faEdit[regular]\:\:Démonstration - #1},
	%title after break=Démonstration (suite),
	colback=white,
	extras={enhanced, frame empty},
	attach boxed title to top left = {xshift = -3mm},
	boxed title style={size=small,colback=violet!65!white,colframe=violet!65!white, arc=3mm,
	}
}

\newtcolorbox{boxPerso}[2]{
	enhanced,
	breakable,
	fonttitle=\bfseries,
	frame empty,
	nobeforeafter,
	borderline west={0.5mm}{0mm}{#2},
	title=#1,
	%title after break=Exemple (suite),
	colback=white,
	extras={enhanced, frame empty},
	detach title,before upper={\tcbtitle\ :\quad },
}


\newcommand*{\Coord}[3]{% 
	\ensuremath{\vv{#1}\, 
		\begin{pmatrix} 
			#2\\ 
			#3 
\end{pmatrix}}}

\newcommand*{\CoordSV}[2]{% 
	\ensuremath{\begin{pmatrix} 
			#1\\ 
			#2 
\end{pmatrix}}}

\newcommand*{\Coordp}[3]{% 
	\ensuremath{\vv{#1}\, 
		\left\lvert 
		\begin{matrix} 
			#2\\ 
			#3 
		\end{matrix}  
		\right.% Ne pas oublier le délimiteur invisible. 
}}


%personnalisation des listes non numérotées pour les différentes boites 
\newlist{itemEx}{itemize}{2}
\setlist[itemEx,1]{label=$\bullet$, font=\color{cyan!75!black}}
\setlist[itemEx,2]{label=$\bullet$, font=\color{cyan!75!black}}

\newlist{enumEx}{enumerate}{2}
\setlist[enumEx,1]{label=\arabic*), font=\bf\color{cyan!75!black}}
\setlist[enumEx,2]{label=\alph*), font=\bf\color{cyan!75!black}}

\newlist{itemMet}{itemize}{2}
\setlist[itemMet,1]{label=$\bullet$, font=\color{blue!75!black}}
\setlist[itemMet,2]{label=$\bullet$, font=\color{blue!75!black}}

\newlist{enumMet}{enumerate}{2}
\setlist[enumMet,1]{label=\arabic*), font=\bf\color{blue!75!black}}
\setlist[enumMet,2]{label=\alph*), font=\bf\color{blue!75!black}}

\newlist{itemDef}{itemize}{2}
\setlist[itemDef,1]{label=$\bullet$, font=\color{ForestGreen}}
\setlist[itemDef,2]{label=$\bullet$, font=\color{ForestGreen}}

\newlist{enumDef}{enumerate}{2}
\setlist[enumDef,1]{label=\arabic*), font=\bf\color{ForestGreen}}
\setlist[enumDef,2]{label=\alph*), font=\bf\color{ForestGreen}}

\newlist{itemProp}{itemize}{2}
\setlist[itemProp,1]{label=$\bullet$, font=\color{red!85!white}}
\setlist[itemProp,2]{label=$\bullet$, font=\color{red!85!white}}

\newlist{enumProp}{enumerate}{2}
\setlist[enumProp,1]{label=\arabic*), font=\bf\color{red!85!white}}
\setlist[enumProp,2]{label=\alph*), font=\bf\color{red!85!white}}

\newlist{itemDem}{itemize}{2}
\setlist[itemDem,1]{label=$\bullet$, font=\color{violet!85!white}}
\setlist[itemDem,2]{label=$\bullet$, font=\color{violet!85!white}}

\newlist{enumDem}{enumerate}{2}
\setlist[enumDem,1]{label=\arabic*), font=\bf\color{violet!85!white}}
\setlist[enumDem,2]{label=\alph*), font=\bf\color{violet!85!white}}

\newlist{itemRem}{itemize}{2}
\setlist[itemRem,1]{label=\ding{43}, font=\color{orange!85!white}}
\setlist[itemRem,2]{label=\ding{43}, font=\color{orange!85!white}}

\newlist{enumRem}{enumerate}{2}
\setlist[enumRem,1]{label=\arabic*), font=\bf\color{orange!85!white}}
\setlist[enumRem,2]{label=\alph*), font=\bf\color{orange!85!white}}

% nouvelle commande pour mettre entre guillemets une expression
\newcommand{\guillemets}[1]{\og \emph{#1} \fg}


\newcommand{\corr}[1]{\fcolorbox{red}{yellow!30}{#1}}

\newtcolorbox{corrBox}{enhanced,breakable,boxrule=1pt,nobeforeafter,boxsep=0pt,left=\fboxsep,right=\fboxsep,colback=yellow!30,colframe=red!75!black}

% commandes pour mettre en couleur (et en gras) dans une expression mathématique.
\newcommand{\mVert}[1]{{\color{green!55!black}\bm{#1}}}
\newcommand{\mVertWB}[1]{{\color{green!55!black}#1}}
\newcommand{\mRouge}[1]{{\color{red!85!white}\bm{#1}}}
\newcommand{\mRougeWB}[1]{{\color{red!85!white}#1}}
\newcommand{\mBleu}[1]{{\color{blue!75!black}\bm{#1}}}
\newcommand{\mBleuWB}[1]{{\color{blue!75!black}#1}}
\newcommand{\mOrange}[1]{{\color{orange!75!white}\bm{#1}}}
\newcommand{\mOrangeWB}[1]{{\color{orange!75!white}#1}}
\newcommand{\mViolet}[1]{{\color{violet!85!white}\bm{#1}}}
\newcommand{\mVioletWB}[1]{{\color{violet!85!white}#1}}

%commandes pour les flèches de distributivité.
\newcommand{\tikzmarkBB}[2][]{ \tikz[remember picture, overlay, baseline=-0.5ex]\node[#1, left](#2){}; }
\newcommand{\flechePropV}[4][]{ \tikz[remember picture, overlay]\draw (#2)edge[bend left=90,min distance=1em, -stealth]node[midway, #1]{#4}(#3); }
\newcommand{\flechePropH}[4][]{ \tikz[remember picture, overlay]\draw (#2)edge[bend left=90, min distance=1em, -stealth]node[midway, #1]{#4}(#3); }
\newcommand{\flechePropHbas}[4][]{ \tikz[remember picture, overlay]\draw (#2)edge[bend left=-90, min distance=1em, -stealth]node[midway, #1]{#4}(#3); }


\lstdefinestyle{algo}{%
	inputencoding=utf8,
	mathescape,
	basicstyle=\footnotesize\ttfamily,
	numberstyle=\scriptsize\ttfamily\color{gray},
	showstringspaces=false,
	tabsize=2,
	numbers=left,
	framexleftmargin=8mm,
	xleftmargin=8mm,
	keepspaces=true,
	frame=lines,
	framerule=1pt,
	rulecolor=\color{green!50!black},
	backgroundcolor=\color{green!5},
	keywords = {Pour, allant, Si, alors, Sinon, Fin},    
	morecomment=[l][commentstyle]{//},
	commentstyle=\color{gray},
	classoffset=0,
	keywordstyle=\color{blue!75},
	classoffset=1,
	morekeywords = {Afficher},
	keywordstyle=\color{green!50!black},
	classoffset=0,
	morestring=[b]",
	stringstyle=\color{red!75},
	literate=
	{é}{{\'e}}{1}%
	{è}{{\`e}}{1}%
	{à}{{\`a}}{1}%
	{â}{{\^a}}{1}%%%
	{ç}{{\c{c}}}{1}%
	{œ}{{\oe}}{1}%
	{ù}{{\`u}}{1}%
	{É}{{\'E}}{1}%
	{È}{{\`E}}{1}%
	{À}{{\`A}}{1}%
	{Ç}{{\c{C}}}{1}%
	{Œ}{{\OE}}{1}%
	{Ê}{{\^E}}{1}%
	{ê}{{\^e}}{1}%
	{î}{{\^i}}{1}%
	{ï}{{\"i}}{1}%%%
	{ô}{{\^o}}{1}%
	{û}{{\^u}}{1}%
	{=}{$\leftarrow$}{1}%
	{==}{$={}$}{1}}


%  --------   Macro pour les exercices et leur numérotation  

\newcounter{exo}

\setcounter{exo}{1}

\newcommand{\exercice}[1]{\vskip 0.2cm\noindent{\bf \ Exercice \theexo {  #1 }}\addtocounter{exo}{1}}

\newcounter{rit}

\setcounter{rit}{1}

\newcommand{\rituel}[1]{\vskip 0.2cm\noindent{ \bf \ \underline{Rit. \therit {  #1 }}}\addtocounter{rit}{1}}

\newcommand{\raisedrule}[2][0em]{\leaders\hbox{\rule[#1]{1pt}{#2}}\hfill}

\newcounter{numexo}

\setcounter{numexo}{1}

\newcommand{\exerciceDS}[1]{\vskip 0.6cm\noindent\rule[+2pt]{1cm}{0.5pt}\ {\large\bf \ Exercice \thenumexo {  #1 }}\raisedrule[+2pt]{0.5pt}\null\addtocounter{numexo}{1}\par\vskip 0,4cm}

\newcommand{\bareme}[1]{\tikz\node[circle,draw,inner sep=1pt,outer sep=1pt, color=red]{\textbf{#1}};}

%  --------   Pour le langage Python

% Default fixed font does not support bold face
\DeclareFixedFont{\ttb}{T1}{txtt}{bx}{n}{12} % for bold
\DeclareFixedFont{\ttm}{T1}{txtt}{m}{n}{12}  % for normal

% Custom colors
\usepackage{color}
\definecolor{deepblue}{rgb}{0,0,0.5}
\definecolor{deepred}{rgb}{0.6,0,0}
\definecolor{deepgreen}{rgb}{0,0.5,0}

\usepackage{listings}

% Python style for highlighting
\newcommand\pythonstyle{\lstset{
		language=Python,
		basicstyle=\ttm,
		otherkeywords={self},             % Add keywords here
		keywordstyle=\ttb\color{deepblue},
		emph={MyClass,__init__},          % Custom highlighting
		emphstyle=\ttb\color{deepred},    % Custom highlighting style
		stringstyle=\color{deepgreen},
		frame=tb,                         % Any extra options here
		showstringspaces=false            % 
}}


% Python environment
\lstnewenvironment{python}[1][]
{
	\pythonstyle
	\lstset{#1}
}
{}

% Python for external files
\newcommand\pythonexternal[2][]{{
		\pythonstyle
		\lstinputlisting[#1]{#2}}}

% Python for inline
\newcommand\pythoninline[1]{{\pythonstyle\lstinline!#1!}}


% \begin{document}
	
	% \section{``In-text'' listing highlighting}
	
	% \begin{python}
		% class MyClass(Yourclass):
		%     def __init__(self, my, yours):
		%         bla = '5 1 2 3 4'
		%         print bla
		% \end{python}
	
	% \section{External listing highlighting}
	
	% \pythonexternal{demo.py}
	
	% \section{Inline highlighting}
	
	% Definition \pythoninline{class MyClass} means \dots
	
	% \end{document}

\definecolor{codegreen}{rgb}{0,0.6,0}
\definecolor{codegray}{rgb}{0.5,0.5,0.5}
\definecolor{codepurple}{rgb}{0.58,0,0.82}
\definecolor{backcolour}{rgb}{0.95,0.95,0.92}

\lstdefinestyle{myStyle}{
	backgroundcolor=\color{backcolour},   
	commentstyle=\color{codegreen},
	keywordstyle=\color{magenta},
	numberstyle=\tiny\color{codegray},
	stringstyle=\color{codepurple},
	basicstyle=\ttfamily\footnotesize,
	breakatwhitespace=false,         
	breaklines=true,                 
	captionpos=b,                    
	keepspaces=true,                 
	numbers=left,                    
	numbersep=5pt,                  
	showspaces=false,                
	showstringspaces=false,
	showtabs=false,                  
	tabsize=2
}

\lstset{style=myStyle}

\newcommand{\normeB}[1]{\left\Vert \vv{#1}\right\Vert}

\newcommand\grasmaths[1]{\symbf{#1}}
\newcommand\sfbmath[1]{\symbf{\symsfup{#1}}}
\definecolor{marron}{rgb}{0.64,0.16,0.16}
\definecolor{titrebleu}{rgb}{0,0,0.6}

\usepackage{tabularx}
\usepackage{wasysym}
% ===========================tableaux de base
\newcolumntype{R}[1]{>{\raggedleft\arraybackslash }b{#1}} 
\newcolumntype{L}[1]{>{\raggedright\arraybackslash }b{#1}}
\newcolumntype{C}[1]{>{\centering\arraybackslash }b{#1}}
\newcolumntype{M}[1]{>{\centering\arraybackslash }m{#1}} %doublecentrée
\renewcommand{\tabularxcolumn}[1]{m{#1}}
\newcolumntype{Y}{>{\centering\arraybackslash }X}

%=============================boîtes de cours
\usepackage{fontawesome5}
\newcommand\fontbox{\sffamily}
%\newcommand\titrebox[1]{\textbf{#1}}
\tcbuselibrary{breakable}
\tcbset{bdefaut/.style={%
		fonttitle=\fontbox,%
		enhanced,breakable,center,boxrule=1.25pt,drop shadow=black!33!white,%
		sharp corners=uphill,arc=12pt,before skip=12pt,after skip=12pt,%
		left=5pt,leftupper=5pt,top=1mm,bottom=1mm,%
		before upper=\leavevmode
	}
}
\NewDocumentCommand\bbnewbox{ m m m m O{#4} !O{} }{%nom + titre + icone + bordure/titre/fondtitre + couleur fond + options
	\DeclareTColorBox{#1}{ O{} }{%
		bdefaut,%
		colframe=#4,coltitle=#4!25!Black,colbacktitle=#4!35,colback=#5!10,%
		title={#3\:\:\titrebox{#2##1}},%
		#6
	}%
}

\bbnewbox{bdefi}{Définition}{\faIcon[regular]{comment-dots}}{ForestGreen}
\bbnewbox{bcor}{Corollaire}{\faIcon[regular]{comment-dots}}{ForestGreen}
\bbnewbox{bdefis}{Définitions}{\faIcon[regular]{comment-dots}}{ForestGreen}
\bbnewbox{bhumour}{Humour}{\faIcon[regular]{laugh-wink}}{Yellow}
%\bbnewbox{bprop}{Propriété}{\faCog}{Salmon}
\bbnewbox{bprop}{Propriété}{\faCog}{red!85!white}
\bbnewbox{bprops}{Propriétés}{\faCog}{Salmon}
\bbnewbox{bcscq}{Conséquence}{\faArrowAltCircleRight[regular]}{Salmon}
\bbnewbox{bthm}{Théorème}{\faBullhorn}{FireBrick}
\bbnewbox{blog}{Logiciel}{\faDev}{CornflowerBlue}[white]
\bbnewbox{bhistoire}{Point histoire}{\faUniversity}{Chocolate}[white]

%\bbnewbox{brmq}{Remarque}{\faHandPointRight[regular]}{MidnightBlue}[white]
\bbnewbox{bexercice}{Exercice}{\faHubspot}{MidnightBlue}
%\bbnewbox{bexemple}{Exemple}{\faHubspot}{MediumTurquoise}
\bbnewbox{bvoc}{Vocabulaire}{\faList*[regular]}{YellowGreen}
\bbnewbox{bnota}{Notation}{\faList*}{YellowGreen}
\bbnewbox{bregle}{Règle}{\faList*}{YellowGreen}
\bbnewbox{bcadre}{Cadre}{\faInfoCircle}{Peru}
\bbnewbox{bidee}{Idée}{\faInfo}{LightBlue}
\bbnewbox{bintro}{Intro}{\faMarker}{LightSalmon}
\bbnewbox{billustr}{Illustration}{\faChartLine}{Crimson}[White]
\bbnewbox{bvideo}{}{\faIcon[regular]{file-video}}{teal}[White]
\bbnewbox{bbilan}{Bilan}{\faForward}{Purple}
%\bbnewbox{bmethode}{Méthode}{\faTasks}{Coral}
%\bbnewbox{bmethode}{Méthode}{\faTasks}{blue!75!black}
\bbnewbox{btuto}{Tuto}{\faTerminal}{Coral}[White]
\bbnewbox{bmanip}{Manipulation}{\faCodeBranch}{DarkCyan}[White]
\bbnewbox{bcalco}{Calculatrice}{\faCalculator}{MediumVioletRed}[White]
\bbnewbox{boutil}{Outil}{\faTools}{DeepSkyBlue}[White]
\bbnewbox{basavoir}{À savoir}{\faHeart}{Red}
\bbnewbox{bbloctd}{TD}{\faHourglassStart}{LightSeaGreen}[White]
\bbnewbox{bblocdm}{À vous de jouer}{\faHourglassStart}{LightSeaGreen}[White]
%\bbnewbox{brappel}{Rappel}{\faHandPointUp[regular]}{LightSalmon}
\bbnewbox{brappel}{Rappel}{\faHandPointUp[regular]}{Peru}
\bbnewbox{bresol}{Résolution}{\faPuzzlePiece}{LightSeaGreen}[White]
\bbnewbox{battention}{Attention}{\faExclamationTriangle}{Red}
\bbnewbox{bcode}{Code}{\faCode}{Crimson}[White]
\bbnewbox{bpython}{Python}{\faPython}{SeaGreen}[White]

\newcommand{\grasInDef}[1]{\textcolor{ForestGreen!55!black}{\textbf{#1}}}
\newcommand{\grasInProp}[1]{\textcolor{red!85!white}{\textbf{#1}}}
\newcommand{\grasInEx}[1]{\textcolor{MediumTurquoise!55!black}{\textbf{#1}}}
\newcommand{\grasInRem}[1]{\textcolor{orange!65!white}{\textbf{#1}}}
\newcommand{\grasInMet}[1]{\textcolor{blue!75!black}{\textbf{#1}}}
\newcommand{\grasInDem}[1]{\textcolor{violet!85!white}{\textbf{#1}}}

%\newtcolorbox{bmethodet}[1][]{%
%	bdefaut,%
%	colframe=blue!75!black,coltitle=white,colbacktitle=blue!65!black,colback=blue!75!black!15!white,
%	title={\faTasks\:\:\titrebox{{Méthode - #1}}}
%}

%outil
\newtcolorbox{boutilcalc}[1][]{%
	bdefaut,%
	colframe=DeepSkyBlue,coltitle=black,colbacktitle=DeepSkyBlue!35,colback=White,
	title=\textcolor{DeepSkyBlue!25!Black}{\faFileVideo[regular]\:\:\titrebox{Tuto#1 Vidéo}}
}

%loglogo
\newtcolorbox{bloglogo}[2][]{%
	bdefaut,%
	colframe=CornflowerBlue,coltitle=black,colbacktitle=CornflowerBlue!35,colback=white,
	title=\textcolor{CornflowerBlue!25!Black}{\faDev\:\:\titrebox{Logiciel#1}},
	overlay={\begin{tcbclipinterior}
			\node[opacity=0.33,rotate=-45] at ($(interior.south east)+(-10pt,10pt)$){\includegraphics[height=24pt]{#2}};
	\end{tcbclipinterior}}
}

%demo
%\newtcolorbox{bdemo}[1][]{%
%	bdefaut,%
%	colframe=Violet,coltitle=black,colbacktitle=Violet!35,colback=Violet!15,
%	title=\textcolor{Violet!25!Black}{\faEdit[regular]\:\:\titrebox{Démonstration#1}},
%	overlay={\begin{tcbclipinterior}%
%			\node[] at ($(interior.south east)+(-10pt,10pt)$){\large\textcolor{Violet}{\textbf{\CheckedBox}}};%
%	\end{tcbclipinterior}}
%}


%\newtcolorbox{bdemo}[1][]{%
%	bdefaut,%
%	colframe=violet!85!white,coltitle=black,colbacktitle=Violet!35,colback=Violet!15,
%	title=\textcolor{Violet!25!Black}{\faEdit[regular]\:\:\titrebox{Démonstration#1}},
%	overlay={\begin{tcbclipinterior}%
%			\node[] at ($(interior.south east)+(-10pt,10pt)$){\large\textcolor{Violet}{\textbf{\CheckedBox}}};%
%	\end{tcbclipinterior}}
%}

%demoblanc
\newtcolorbox{bdemoblanc}[1][]{%
	bdefaut,%
	colframe=Violet,coltitle=black,colbacktitle=Violet!35,colback=White,
	title=\textcolor{Violet!25!Black}{\faEdit[regular]\:\:\titrebox{Démonstration#1}},
	overlay={\begin{tcbclipinterior}%
			\node[] at ($(interior.south east)+(-10pt,10pt)$){\large\textcolor{Violet}{\textbf{\CheckedBox}}};%
	\end{tcbclipinterior}}
}

%algo
\newtcolorbox{balgo}[1][]{%
	bdefaut,%
	colframe=OrangeRed,coltitle=black,colbacktitle=OrangeRed!35,colback=White,
	title=\textcolor{OrangeRed!25!Black}{\faFileCode\:\:\titrebox{Algorithme#1}},
	overlay={\begin{tcbclipinterior}%
			\node[opacity=0.15,rotate=-15] at ($(interior.south east)+(-8pt,12pt)$){\includegraphics[height=24pt]{algorithmvec}};%
	\end{tcbclipinterior}}
}


%aide
\newtcolorbox{baide}{%
	bdefaut,%
	colframe=titrebleu!35,coltitle=titrebleu!50!Black,colbacktitle=titrebleu!15,colback=White,
	title=\textcolor{titrebleu!25!Black}{\faBacon\:\:\titrebox{Niveaux de difficulté}},
	overlay={\begin{tcbclipinterior}%
			\node[opacity=0.33,rotate=45,scale=0.33] at ($(interior.south east)+(-6pt,6pt)$){\textcolor{titrebleu!50!Black}{\footnotesize \bfseries \sffamily \textcopyright{} Dairon11}};%
	\end{tcbclipinterior}}
}

\newtcolorbox{bcalcologo}[2][]{%
	bdefaut,%
	colframe=MediumVioletRed,coltitle=black,colbacktitle=MediumVioletRed!35,colback=White,
	title=\textcolor{MediumVioletRed!25!Black}{\faCalculator\:\:\titrebox{Calculatrice#1}},
	overlay={\begin{tcbclipinterior}%
			\node[opacity=0.33,rotate=-45] at ($(interior.south east)+(-10pt,10pt)$){\includegraphics[height=48pt]{#2}};%
	\end{tcbclipinterior}}
}

\definecolor{orangescratch}{HTML}{F8A31F}
\newtcolorbox{bscratch}[1][]{%
	bdefaut,%
	colframe=orangescratch,coltitle=black,colbacktitle=orangescratch!35,colback=White,
	title=\textcolor{orange!25!Black}{\includegraphics[height=1em]{scratch}\:\:\titrebox{Scratch#1}},
	overlay={\begin{tcbclipinterior}%
			\node[opacity=0.33,rotate=45] at ($(interior.south east)+(-12pt,12pt)$){\includegraphics[height=12pt]{scratchban}};%
	\end{tcbclipinterior}}
}

%reponse
\newtcolorbox{breponse}[1][]{%
	bdefaut,%
	colframe=Thistle,coltitle=black,colbacktitle=Thistle!35,colback=White,
	before upper = \vspace{-0.05cm},
	title=\textcolor{Thistle!25!Black}{\faIcon{pen-alt}\:\:\titrebox{Réponse#1}}
}

\usepackage{dashundergaps}
\newcommand\underdash[1][black]{% <- règle la couleur du texte
	\bgroup
	\ifdim\ULdepth=\maxdimen\settodepth\ULdepth{(j}\advance\ULdepth.4pt\fi
	\markoverwith{\kern0.2em% <- règle l'espacement avant un pointillé
		\vtop{\kern0.5ex% <- règle l'altitude des pointillés
			{\color{black}% <- règle la couleur des pointillés
				\hrule width.4em% <- règle la largeur des pointillés
		}}\kern0.2em% <- règle l'espacement après un pointillé
	}\color{#1}\ULon}

\newcommand\pointille[1][1cm]{\makebox[#1]{\dotfill}}


\newboolean{versionprof}
\setboolean{versionprof}{true}

\usepackage{epic}   % extension et simplification de  l'environnement picture
\usepackage{graphicx}    % permet d'inclure des graphique externe dans un document
\usepackage{arydshln}   %permet d'avoir des lignes en pointillés

\newcommand{\lignes}[1]{
	\setlength{\unitlength}{1mm}
	\vskip 0,4cm\noindent
	\multido{\i=0+1}{#1}{%
		\noindent
		\begin{picture}(150,5)
			\linethickness{0,005mm}
			\put(-5,0){\dashline{1}(175,5)(5,5)}
		\end{picture}\vskip 0,001cm
	}
	\vskip 0cm
}

\rhead{Rappels de cours sur les vecteurs}

\newcommand{\va}[1]{\lvert \ #1 \ \rvert}
\newcommand{\vaf}[1]{\left\lvert \ #1\  \right\rvert}

\begin{document}
	\graphicspath{{../images/}}
	
	\begin{center}
		
		
		\fbox{
			\begin{minipage}[]{15cm}
				\begin{center}
					\LARGE{\rule{0.5cm}{0pt}\rule[-0.25cm]{0pt}{0.9cm}
						Vecteurs\\
						Rappels de cours
					}
				\end{center}
			\end{minipage}
		}
	\end{center}

\section{géométrie dans un repère}


% \subsection{Repères}



% \begin{definition}{(Repère du plan)}
% 	\noindent Trois points \textbf{distincts} et \textbf{non alignés} O, I et J forment un {\ifthenelse{\boolean{versionprof}}{\emph{repère du plan}}{\quad\dots\phantom{repère du plan xxxxxxx}}} que l'on note {\ifthenelse{\boolean{versionprof}}{(O,I,J)}{\ \dots\phantom{(O,I,J) xxx}}}.
	
% 	\vspace{0.2cm}
% 	\noindent Les unités OI et OJ permettent de graduer les axes (OI) et (OJ).
	
% 	\begin{itemDef}
% 		\vspace{0.1cm}
% 		\item Le point O s'appelle {\ifthenelse{\boolean{versionprof}}{\textbf{l'origine}}{$\quad\dots\dots\ $ \hspace{2.cm}}} du repère.
% 		\vspace{0.1cm}
		
% 		\item L'axe $(OI)$ s'appelle l'axe des {\ifthenelse{\boolean{versionprof}}{\textbf{abscisses}}{$\quad\dots\dots\ $ \hspace{2.cm}}} et la distance OI donne l'unité sur cet axe.
% 		\vspace{0.1cm}
% 		\item L'axe $(OJ)$ s'appelle l'axe des {\ifthenelse{\boolean{versionprof}}{\textbf{ordonnées}}{$\quad\dots\dots\ $ \hspace{2.cm}}} et la distance OJ donne l'unité sur cet axe.
% 	\end{itemDef}
% \end{definition}


% \begin{center}
% 	\psset{unit=.5cm}
% 	\begin{pspicture}(-2,-3)(6,6)
% 		\psline[linewidth=.75pt,linecolor=darkgray]{->}(-2,0)(6,0)
% 		\psline[linewidth=.75pt,linecolor=darkgray]{->}(-.5,-1)(3,6)
		
% 		\psdots[dotstyle=*](1.5,0)(3,0)(4.5,0)
% 		\psdots[dotstyle=*](1,2)(2,4)
		
% 		\psline[linewidth=1.5pt,linecolor=black]{-}(0,0)(1.5,0)
% 		\psline[linewidth=1.5pt,linecolor=black]{-}(0,0)(1,2)
		
% 		\begin{scriptsize}
% 			\uput[d](1.5,0){I}
% 			\uput[l](1,2){J}
% 			\uput[dl](-0.1,0){O}
% 			\uput[dl](6,0){$x$}
% 			\uput[dl](2.8,6){$y$}
% 		\end{scriptsize}
% 		{\ifthenelse{\boolean{versionprof}}{\rput(2,-2.5){Repère \textbf{quelconque}}}{\rput(2,-2.5){Repère \quad\dots\phantom{quelconque}}}}
% 	\end{pspicture} \hfill
% 	\psset{unit=.5cm}
% 	\begin{pspicture}(-2,-3)(6,6)
% 		\psaxes[labels=none,linewidth=.75pt,ticks=none,linecolor=darkgray]{->}(0,0)(-2,-1)(6,6)
% 		\psaxes[labels=none,linewidth=1.5pt,ticks=none,linecolor=black]{-}(0,0)(1.5,2.6)
% 		\psdots[dotstyle=*](1.5,0)(3,0)(4.5,0)
% 		\psdots[dotstyle=*](0,2.6)(0,5.2)
		
% 		\begin{scriptsize}
% 			\uput[dl](0,0){O}
% 			\uput[dl](6,0){$x$}
% 			\uput[dl](0,6){$y$}
% 			\uput[d](1.5,0){I}
% 			\uput[l](0,2.6){J}
% 		\end{scriptsize}
		
% 		{\ifthenelse{\boolean{versionprof}}{\rput(2,-2.5){Repère \textbf{orthogonal}}}{\rput(2,-2.5){Repère \quad\dots\phantom{orthogonal}}}}
% 		{\ifthenelse{\boolean{versionprof}}{\rput(2,-3.5){(OI) $\perp$ (OJ)}}{\rput(2,-3.5){\dots}}}
% 	\end{pspicture}\hfill
% 	\psset{unit=.5cm}
% 	\begin{pspicture}(-2,-3)(6,6)
% 		\psaxes[labels=none,linewidth=.75pt,ticks=none,linecolor=darkgray]{->}(0,0)(-2,-1)(6,6)
% 		\psaxes[labels=none,linewidth=1.5pt,ticks=none,linecolor=black]{-}(0,0)(2,2)
% 		\psdots[dotstyle=*](2,0)(4,0)
% 		\psdots[dotstyle=*](0,2)(0,4)
		
% 		\begin{scriptsize}
% 			\uput[dl](0,0){O}
% 			\uput[dl](6,0){$x$}
% 			\uput[dl](0,6){$y$}
% 			\uput[d](2,0){I}
% 			\uput[l](0,2){J}
% 		\end{scriptsize}
		
% 		{\ifthenelse{\boolean{versionprof}}{\rput(2,-2.5){Repère \textbf{orthonormé}}}{\rput(2,-2.5){Repère \quad\dots\phantom{orthonormé}}}}
% 		{\ifthenelse{\boolean{versionprof}}{\rput(2,-3.5){(OI) $\perp$ (OJ) \quad et \quad OI $=$ OJ}}{\rput(2,-3.5){\dots\quad et \quad\dots}}}
% 	\end{pspicture}
% \end{center}


\subsection{Milieu d'un segment}
\begin{bprop}
	\begin{minipage}[c]{.65\linewidth}
		
		Dans un repère \textbf{quelconque} du plan :
		Si $A$ et $B$ sont deux points de coordonnées respectives $(x_A;y_A)$ et $(x_B;y_B)$ :
		
		\medskip
		\noindent le milieu de $[AB]$ a pour coordonnées  {\ifthenelse{\boolean{versionprof}}{$\Big(\ \dfrac{x_A\ + \ x_B}{2}\ ; \ \dfrac{y_A\ + \ y_B}{2}\ \Big)$.}{$\Big(\ \dfrac{\phantom{xxx + xxx}}{\phantom{2}} \ ; \ \dfrac{\phantom{xxx + xxx}}{\phantom{2}}\ \Big)$.}}  
		
	\end{minipage}\hfill
	% image à droite - surtout pas de ligne vide avant cette partie        
	\begin{minipage}[c]{.35\linewidth}
		\vspace{-0.5cm}
		\begin{center}
			\psset{xunit=0.65cm,yunit=0.65cm,algebraic=true,dimen=middle,dotstyle=o,dotsize=5pt 0,linewidth=1.pt,arrowsize=3pt 2,arrowinset=0.25}
			\begin{pspicture*}(-1,-0.8)(5.5,6.5)
				
				% pas besoin de grille, ça encombre
				%\psgrid[gridwidth=0.8pt,gridcolor=darkgray,subgriddiv=0,gridlabels=0,griddots=4](0,0)(-2.,-1.)(5.,5.)
				
				\psaxes[labelFontSize=\scriptstyle,xAxis=true,yAxis=true,labels=none,Dx=1.,Dy=1.,ticksize=-4pt 0,subticks=0]{->}(0,0)(-0.5,-0.5)(5.5,6)
				
				\psline[linewidth=1.2pt,linestyle=dotted](2,0)(2,2)
				\psline[linewidth=1.2pt,linestyle=dotted](4,0)(4,5)
				\psline[linewidth=1.2pt,linestyle=dotted](0,2)(2.,2)
				\psline[linewidth=1.2pt,linestyle=dotted](0,5)(4.,5)
				\psline[linewidth=1.2pt,linestyle=dotted](2,2)(4.,5)
				\psline[linewidth=1.2pt,linestyle=dotted](0,3.5)(3.,3.5)
				\psline[linewidth=1.2pt,linestyle=dotted](3,0)(3.,3.5)
				
				\psdots[dotsize=3pt 0,dotstyle=*](0.,0.)(1.,0.)(0.,1.)(2.,2.)(4.,5)
				
				\begin{scriptsize}
					\uput[dl](0,0){O}
					\uput[d](1,0){I}
					\uput[l](0,1){J}
					\uput[ul](2.,2){A}
					\uput[dr](4,5){B}
					\uput[dl](5.5,0){$x$}
					\uput[d](2,0){$x_A$}
					\uput[d](4,0){$x_B$}
					\uput[l](0,2){$y_A$}
					\uput[l](0,5){$y_B$}
					\uput[dl](0,6){$y$}
				\end{scriptsize}
			\end{pspicture*}
		\end{center}
	\end{minipage}
\end{bprop}

\bigskip

\begin{bmethodet}{Milieu et parallélogramme}
	Rappel : les diagonales d'un parallélogramme {\ifthenelse{\boolean{versionprof}}{se coupent en leur milieu.}{\quad \dots}}
	
	\medskip
	Pour montrer qu'un quadrilatère ABCD est un parallélogramme, on montre que les 2 diagonales [AC] et [BD] ont le même milieu :
	\begin{enumMet}
		\item On détermine les coordonnées du milieu M de la diagonale [AC].
		\item On détermine les coordonnées du milieu N de la diagonale [BD].
		\item On montre que les points M et N sont confondus et on conclut.
	\end{enumMet}
	
	Pour trouver le 4ème point D d'un parallélogramme ABCD :
	\begin{enumMet}
		\item On détermine le milieu M de la diagonale [AC].
		\item On détermine le point D tel que M soit milieu de la diagonale [BD].
	\end{enumMet}
	
\end{bmethodet}


\begin{bexemple} : Dans le repère (O;I;J), on veut déterminer les coordonnées de M milieu de [AB] et N milieu de [AC].
	
	\begin{minipage}[c]{.7\linewidth}
		Les coordonnées de A, B et C sont : A$(\ 1 \ ; \ 4 \ )$ \quad B$(\ 5 \ ; \ 2 \ )$ \quad C$(\ -2 \ ; \ -2 \ )$
		
		\medskip
		Les coordonnées de M milieu de [AB] sont :
		
		$x_M=$ {\ifthenelse{\boolean{versionprof}}{
				$\dfrac{x_A\ + \ x_B}{2}=\dfrac{1\ + \ 5}{2}=\dfrac{6}{2}=3$}
			{$\dfrac{x_A\ + \ x_B}{2}=\dfrac{\phantom{xxx}+\phantom{xxx}}{\phantom{2}}=\dfrac{\phantom{xxx}}{\phantom{2}}=\dots$}} 
		
		$y_M=$ {\ifthenelse{\boolean{versionprof}}{
				$\dfrac{y_A\ + \ y_B}{2}=\dfrac{4\ + \ 2}{2}=\dfrac{6}{2}=3$ \quad donc M$(\ 3 \ ; \ 3 \ )$.}
			{$\dfrac{y_A\ + \ y_B}{2}=\dfrac{\phantom{xxx}+\phantom{xxx}}{\phantom{2}2}=\dfrac{\phantom{xxx}}{\phantom{2}}=\dots$}}
		
		\medskip
		Les coordonnées de N milieu de [AC] sont :
		
		$x_N=$ {\ifthenelse{\boolean{versionprof}}{
				$\dfrac{x_A\ + \ x_C}{2}=\dfrac{1\ + \ (-2)}{2}=-\dfrac{1}{2}=-0,5$}
			{$\dfrac{x_A\ + \ x_C}{2}=\dfrac{\phantom{xxx}+\phantom{xxx}}{\phantom{2}}=\dfrac{\phantom{xxx}}{\phantom{2}}=\dots$}}
		
		$y_N=$ {\ifthenelse{\boolean{versionprof}}{
				$\dfrac{y_A\ + \ y_C}{2}=\dfrac{4\ + \ (-2)}{2}=\dfrac{2}{2}=1$ \quad donc N$(\ -0,5 \ ; \ 1 \ )$.}
			{$\dfrac{y_A\ + \ y_C}{2}=\dfrac{\phantom{xxx}+\phantom{xxx}}{\phantom{2}2}=\dfrac{\phantom{xxx}}{\phantom{2}}=\dots$}}
		
		\medskip
		On peut le vérifier graphiquement en traçant les segments et leur milieu.
	\end{minipage}
	\hfill
	\begin{minipage}[c]{.3\linewidth}
		\begin{center}
			\psset{xunit=0.6cm,yunit=0.6cm,algebraic=true,dimen=middle,dotstyle=o,dotsize=5pt 0,linewidth=1.pt,arrowsize=3pt 2,arrowinset=0.25}
			\begin{pspicture*}(-3.,-2.5)(6,5)
				
				\psgrid[gridwidth=0.8pt,gridcolor=darkgray,subgriddiv=0,gridlabels=0,griddots=4](0,0)(-3.,-2.)(6.,5.)
				
				\psaxes[labelFontSize=\scriptstyle,xAxis=true,yAxis=true,labels=none,Dx=1.,Dy=1.,ticksize=-4pt 0,subticks=2]{->}(0,0)(-3.,-2)(6,5)
				
				\psdots[dotsize=3pt 0,dotstyle=*](0.,0.)(1.,0.)(0.,1.)(1.,4.)(5,2)(-2,-2)
				
				{\ifthenelse{\boolean{versionprof}}{
						\pspolygon[linecolor=red](1.,4.)(5,2)(-2,-2)}{}}  
				
				\begin{scriptsize}
					\uput[dl](0,0){O}
					\uput[d](1,0){I}
					\uput[l](0,1){J}
					\uput[ur](1,4){A}
					\uput[ur](5,2){B}
					\uput[ul](-2,-2){C}
					\uput[dl](5,0){$x$}
					\uput[dl](0,5){$y$}
				\end{scriptsize}
			\end{pspicture*}
		\end{center}
	\end{minipage}
\end{bexemple}

\subsection{Distance entre deux points dans un repère orthonormé}
\begin{bprop}
	\bigskip
	\begin{minipage}[c]{.7\linewidth}
		
		\noindent Soit $A$ et $B$ deux points de coordonnées respectives $(x_A;y_A)$ et $(x_B;y_B)$ dans un \textbf{repère orthonormé} :
		
		\begin{center} $AB \quad = \quad {\ifthenelse{\boolean{versionprof}}{ \sqrt{\quad \left( \ x_B \ - \ x_A \ \right)^2  \quad + \quad \left( \ y_B \ - \ y_A \ \right)^2 \  }}{ \sqrt{\quad \big(\ \phantom{xxxx \ - \ xxxx} \ \big) \phantom{^2  \quad + \quad} \big( \ \phantom{xxxx  \ - \ xxxx} \ \big)\phantom{^2} \quad }}} \quad$\end{center}
		
		\smallskip
		{\small Attention : cette formule est \\
			\underline{valable uniquement si le repère est orthonormé}.}
		
	\end{minipage}\hfill
	% image à droite - surtout pas de ligne vide avant cette partie        
	\begin{minipage}[c]{.3\linewidth}
		\vspace{-0.8cm}
		\begin{center}
			\psset{xunit=0.65cm,yunit=0.65cm,algebraic=true,dimen=middle,dotstyle=o,dotsize=5pt 0,linewidth=1.pt,arrowsize=3pt 2,arrowinset=0.25}
			\begin{pspicture*}(-1,-0.8)(6.5,6.5)
				
				% pas besoin de grille, ça encombre
				%\psgrid[gridwidth=0.8pt,gridcolor=darkgray,subgriddiv=0,gridlabels=0,griddots=4](0,0)(-2.,-1.)(5.,5.)
				
				\psaxes[labelFontSize=\scriptstyle,xAxis=true,yAxis=true,labels=none,Dx=1.,Dy=1.,ticksize=-4pt 0,subticks=0]{->}(0,0)(-0.5,-0.5)(5.5,6)
				
				\psline[linewidth=1.2pt,linestyle=dotted](2,0)(2,5)
				\psline[linewidth=1.2pt,linestyle=dotted](4,0)(4,5)
				\psline[linewidth=1.2pt,linestyle=dotted](0,2)(4.,2)
				\psline[linewidth=1.2pt,linestyle=dotted](0,5)(4.,5)
				\psline[linewidth=1.2pt,linestyle=dotted](2,2)(4.,5)
				
				\psdots[dotsize=3pt 0,dotstyle=*](0.,0.)(1.,0.)(0.,1.)(2.,2.)(4.,5)
				{\ifthenelse{\boolean{versionprof}}{
						\psdots[dotsize=3pt 0,dotstyle=*](4,2)
						\psline[linewidth=0.8pt,linecolor=red]{<->}(2,1)(4,1)
						\psline[linewidth=0.8pt,linecolor=red]{<->}(4.5,2)(4.5,5)
					}{}} 
				
				\begin{scriptsize}
					\uput[dl](0,0){O}
					\uput[d](1,0){I}
					\uput[l](0,1){J}
					\uput[dl](2.,2){A}
					\uput[dr](4,5){B}
					\uput[dl](5.5,0){$x$}
					\uput[d](2,0){$x_A$}
					\uput[d](4,0){$x_B$}
					\uput[l](0,2){$y_A$}
					\uput[l](0,5){$y_B$}
					\uput[dl](0,6){$y$}
					
					{\ifthenelse{\boolean{versionprof}}{
							\uput[dr](4,2){\red H}
							\uput[d](3,1){\red $x_B-x_A$}
							\uput[r](4.3,3.5){\red $y_B-y_A$}
						}{}} 
					
				\end{scriptsize}
			\end{pspicture*}
		\end{center}
	\end{minipage}
	
\end{bprop}

\begin{bexemple} : Les coordonnées de A, B et C sont :
	A$(\ 1 \ ; \ 4 \ )$ \quad B$(\ 5 \ ; \ 2 \ )$ \quad C$(\ -2 \ ; \ -2 \ )$.\\
	\begin{itemEx}
		\item Calculons les longueurs AB, AC et BC.
		\medskip
		Le repère (O,I,J) ci-dessus est \underline{orthonormé} donc on peut utiliser la formule ci-dessus :
		
		\medskip
		AB {\ifthenelse{\boolean{versionprof}}{
				$= \sqrt{ \left(x_B -x_A\right)^2   +  \left(y_B-y_A\right)^2}= \sqrt{ \left(5-1\right)^2   +  \left(2-4\right)^2}= \sqrt{4^2 + (-2)^2 }=\sqrt{16+4 }=\sqrt{20}=2\sqrt{5} $}
			{$= \sqrt{ \left(\phantom{xxx} - \phantom{xxx} \right)^2   +  \left(\phantom{xxx} -\phantom{xxx} \right)^2}= \sqrt{ \left(\phantom{xxx} - \phantom{xxx} \right)^2   +  \left(\phantom{xxx} -\phantom{xxx} \right)^2}= \sqrt{\phantom{xxx} +\phantom{xxx}}=\sqrt{\phantom{xxx} }=\phantom{xxx} $}} 
		
		
		
		
		
		\medskip
		AC {\ifthenelse{\boolean{versionprof}}{
				\textcolor{red}{$= \sqrt{ \left(x_C -x_A\right)^2   +  \left(y_C-y_A\right)^2}= \sqrt{ \left(-2-1\right)^2   +  \left(-2-4\right)^2}= \sqrt{(-3)^2 + (-6)^2 }=\sqrt{9+36 }=\sqrt{45}=3\sqrt{5} $}}
			{$= \sqrt{ \left(\phantom{xxx} - \phantom{xxx} \right)^2   +  \left(\phantom{xxx} -\phantom{xxx} \right)^2}= \sqrt{ \left(\phantom{xxx} - \phantom{xxx} \right)^2   +  \left(\phantom{xxx} -\phantom{xxx} \right)^2}= \sqrt{\phantom{xxx} +\phantom{xxx}}=\sqrt{\phantom{xxx} }=\phantom{xxx} $}}
		
		\medskip
		BC {\ifthenelse{\boolean{versionprof}}{
				\textcolor{red}{$= \sqrt{ \left(x_C -x_B\right)^2   +  \left(y_C-y_B\right)^2}= \sqrt{ \left(-2-5\right)^2   +  \left(-2-2\right)^2}= \sqrt{(-7)^2 + (-4)^2 }=\sqrt{49+16 }=\sqrt{65}$}}
			{$= \sqrt{ \left(\phantom{xxx} - \phantom{xxx} \right)^2   +  \left(\phantom{xxx} -\phantom{xxx} \right)^2}= \sqrt{ \left(\phantom{xxx} - \phantom{xxx} \right)^2   +  \left(\phantom{xxx} -\phantom{xxx} \right)^2}= \sqrt{\phantom{xxx} +\phantom{xxx}}=\sqrt{\phantom{xxx} }=\phantom{xxx} $}}
		
		\item Que peut-on dire du triangle ABC ?\\
		{\ifthenelse{\boolean{versionprof}}{\textcolor{red}{Il semble que ABC soit un triangle rectangle en A. \textbf{Prouvons-le :}}\\
				
				\textcolor{red}{$AB^2+AC^2=(\sqrt{20})^2+(\sqrt{45})^2=20+45=65$ \quad et \quad $BC^2=(\sqrt{65})^2=65$.} 
				
				\textcolor{red}{Donc $AB^2+AC^2=BC^2$ d'après la réciproque du théorème de Pythagore, ABC est rectangle en A}}{}}  
	\end{itemEx} 
\end{bexemple}
\section{Généralités sur les vecteurs}

\subsection{Vecteur et translation}

\begin{bdefi}
	Un vecteur $\vv{AB}$ est un outil mathématique qui matérialise la translation d'un point à un autre point. Il est défini par 3 caractéristiques:
	\begin{itemDef}
		\item sa direction (son inclinaison)
		\item son sens (droite, gauche, haut, bas etc)
		\item sa norme (sa longueur) notée $\| \vv{AB}\|$
	\end{itemDef}
\end{bdefi}

\begin{brmq}
	\grasInRem{Vecteur nul :} le vecteur qui matérialise la translation du point $\mathrm{A}$ au point $\mathrm{A}$, c'està-dire sur lui-même, est noté $\vv{0}$ et s'appelle le vecteur nul. On a $\vv{A A}=\vv{0}$. Ce vecteur n'a ni direction ni sens, et a pour norme $0$.
\end{brmq}
\begin{bdefi}
	\grasInDef{vecteurs égaux : } Deux vecteurs sont égaux s'ils ont : 
	\begin{itemDef}
		\item la même direction
		\item le même sens
		\item la même norme 
	\end{itemDef}
\end{bdefi}

\begin{bdefi}
    \grasInDef{vecteurs opposés : } Deux vecteurs $\vv{u}$ et $\vv{v}$ sont dits opposés s'ils ont :
	\begin{itemDef}
		\item La même direction
		\item La même norme
	\end{itemDef}
	MAIS des sens opposés On note $\vv{u}=-\vv{v}$ ou $\vv{v}=-\vv{u}$
\end{bdefi}
\begin{bprop}
	Le vecteur opposé à $\vv{A B}$ est le vecteur $\vv{B A}$ : c'est-à-dire qu'il part de B pour aller vers A. On a donc:
	$\vv{A B}=-\vv{B A}$ ou encore $\vv{B A}=-\vv{A B}$
\end{bprop}

\section{Opérations sur les vecteurs}
\subsection{Règle du parallélogramme}

\begin{bprop}
	Propriété : ABCD est un parallélogramme équivaut à dire que $\vv{A B}=\vv{D C}$
\end{bprop}

\subsection{Somme de vecteurs et relation de Chasles}

\begin{bdefi}
	\grasInDef{somme de vecteurs : } Soient deux vecteurs $\vv{u}$ et $\vv{v}$ du plan. On définit la somme $\vv{u}+\vv{v}$ comme un vecteur formé en mettant bout à bout les vecteurs $\vv{u}$ et $\vv{v}$.
\end{bdefi}

\begin{bprop}
	\grasInProp{Relation de Chasles : } Soient $A$, $B$ et $C$ trois points du plan. On a : $\vv{A B}+\vv{B C}=\vv{A C}$
\end{bprop}

\begin{bexemple}
	\begin{multicols}{2}
		\begin{enumEx}
			\item $\vv{C B}+\vv{B A}=$ \textcolor{red}{$\vv{CA}$}
			\item $\vv{C B}+\vv{B C}=$ \textcolor{red}{$\vv{CC}=\vv{0}$}
			\item $\vv{A C}+$ \textcolor{red}{$\vv{CB}$} $=\vv{A B}$
			\item $\vv{B C}+\vv{A B}=$ \textcolor{red}{$\vv{AB}+\vv{BC}=\vv{AC}$}
			\item $\vv{C B}-\vv{A B}=$ \textcolor{red}{$\vv{CB}+\vv{BA}=\vv{CA}$}
			\item $\vv{A B}+\vv{\mRougeWB{C}A}=\mRougeWB{\vv{CA}+\vv{AB}=}\vv{C\mRougeWB{B}}$
		\end{enumEx}
	\end{multicols}
\end{bexemple}

\subsection{Multiplication par un réel}

\begin{bprop}
	On peut factoriser ou développer par un réel avec des vecteurs.
\end{bprop}

\begin{bexemples}
	\begin{enumEx}
		\item $4(\vv{u}+\vv{v})=\mRougeWB{4\vv{u}+4\vv{v}}$
		\item $4(2 \vv{u})=\mRougeWB{8\vv{u}}$
		\item $5(2 \vv{u}+3 \vv{v})=\mRougeWB{10\vv{u}+15\vv{v}}$
	\end{enumEx}
\end{bexemples}


\subsection{colinéarité de vecteurs}

\begin{bdefi}
	\grasInDef{notion de colinéarité : } Deux vecteurs $\vv{u}$ et $\vv{v}$ sont dits colinéaires s'ils ont la même direction quel que soit leurs sens et leurs normes. Dans ce cas, il existe un nombre réel $k$ tel que $\vv{u}=k \times \vv{v}$
\end{bdefi}

\begin{bprop}
	\grasInProp{Parallélisme et alignement :} Soient les points $\mathrm{A}, \mathrm{B} \mathrm{C}$ et $\mathrm{D}$ du plan.
	\begin{itemDef}
		\item $\vv{A B}$ et $\vv{C D}$ sont colinéaires équivaut à dire que $(\mathrm{AB}) / /(\mathrm{CD})$
		\item $\vv{A B}$ et $\vv{A C}$ sont colinéaires équivaut à dire que $A, B$ et $C$ sont alignés
	\end{itemDef}
\end{bprop}

\begin{bprop}
	\grasInProp{milieux :} Soient A et B deux points du plan. I est le milieu de $[\mathrm{AB}] \Leftrightarrow \vv{A I}=\dfrac{1}{2} \vv{A B},$ c'est-à-dire $\vv{A I}=\vv{I B}$
\end{bprop}

\section{Coordonnées de vecteurs dans un repère}

\begin{bdefi}
	Soit \Rij, un repère du plan. Pour tout vecteur $\vv{u}$, il existe un unique point $\mathrm{M}$ tel que $\vv{u}=\vv{O M}$. Les coordonnées du vecteur $\vv{u}$ dans ce repère sont alors celles du point $\mathrm{M}$.\\
	Si un point $\mathrm{M}$ a pour coordonnées $\left(x_{M} ; y_{M}\right)$ alors on note  $\vv{u}\left(x_{M} ; y_{M}\right)$ ou $\Coord{u}{x_M}{y_M}$
\end{bdefi}

\begin{center}
	\psset{xunit=0.7cm,yunit=0.7cm,algebraic=true,dimen=middle,dotstyle=o,dotsize=7pt 0,linewidth=2.pt,arrowsize=2pt 2,arrowinset=0.25}
	\begin{pspicture*}(-1.2,-1)(4,3.5)
		%\psgrid[xunit=0.9cm,yunit=0.9cm,subgriddiv=0,gridcolor=lightgray,gridlabels=0pt](-4,-4)(5,5)
		\psaxes[linewidth=0.6pt,labels=none,subticks=1]{->}(0,0)(-1.2,-1)(4,3.5)
		\psline[linewidth=1.2pt,linecolor=blue]{->}(0.,0.)(1.,0.)
		\psline[linewidth=1.2pt,linecolor=blue]{->}(0.,0.)(0.,1.)
		\psline[linewidth=1.2pt,linecolor=red]{->}(0.,0.)(3.,2.)
		\rput[tl](1.8,1.1){\large{\textcolor{red}{$\vv{OM}$}}}
		\psline[linewidth=1.2pt,linecolor=red]{->}(-1.,1.)(2.,3.)
		\rput[tl](0.6,2){\large{\textcolor{red}{$\vv{u}$}}}
		\psdots[dotstyle=x,linecolor=blue](3.,2.)
		\rput[bl](3.1,2.1){\large{\textcolor{blue}{M}}}
		\begin{scriptsize}
			\rput[tl](0.4,-0.1){\textcolor{blue}{$\vv{i}$}}
			\rput[br](-0.1,0.3){\textcolor{blue}{$\vv{j}$}}
			\rput[tr](4,-0.2){\large{$x$}}
			\rput[tl](-0.35,3.5){\large{$y$}}
			\rput[tr](-0.2,-0.2){\large{$O$}}
		\end{scriptsize}
	\end{pspicture*}
\end{center}

\begin{bprop}{égalité de vecteurs}
	Soit $\Coord{u}{x_1}{y_1}$ et $\Coord{v}{x_2}{y_2}$ 2 vecteurs. Alors $: \vv{u}=\vv{v} \iff x_{1}=x_{2}$ et $y_{1}=y_{2}$
\end{bprop}

%\vspace*{-0.4cm}
\subsection{Opérations sur les vecteurs}

\begin{bprop}
	Soit $\Coord{u}{x_1}{y_1}$ et $\Coord{v}{x_2}{y_2}$ 2 vecteurs. Alors:
	\begin{itemProp}
		\item $\vv{u}+\vv{v}\left(\begin{array}{l}
			x_{1}+x_{2} \\
			y_{1}+y_{2}
		\end{array}\right)$
		\item $k \vv{u}\left(\begin{array}{l}
			k\ x_{1} \\
			k\ y_{1}
		\end{array}\right)$
	\end{itemProp}
\end{bprop}

\begin{bmethodet}{Utiliser les opérations sur les vecteurs}
	\grasInMet{Enoncé :}\\
	Soient $\Coord{u}{2}{-7}$ et $\Coord{v}{-3}{4}$. Déterminer les coordonnées du vecteur $\vv{w}=3 \vv{u}-2 \vv{v}$\\
	\grasInMet{Solution :}\\
	Les coordonnées de $\vv{w}$ sont $: 3 \times\left(\begin{array}{c}2 \\ -7\end{array}\right)-2 \times\left(\begin{array}{c}-3 \\ 4\end{array}\right)$\\
	$\Coord{w}{6-(-6)}{-14-8}$\\
	$\Coord{w}{12}{-22}$
\end{bmethodet}

\subsection{Calcul des coordonnées d'un vecteur}

\begin{bprop}
	Dans un repère, soient $\mathrm{A}$ et $\mathrm{B}$ deux points de coordonnées $\mathrm{A}\left(\mathrm{x}_{\mathrm{A}} ; \mathrm{y}_{\mathrm{A}}\right)$ et $\mathrm{B}\left(\mathrm{x}_{\mathrm{B}} ; \mathrm{y}_{\mathrm{B}}\right)$.\\
	alors $\vv{A B}$ a pour coordonnées $\Coord{AB}{x_B-x_A}{y_B-y_A}$.
\end{bprop}

\begin{bmethodet}{Coordonnées d'un vecteur}
	\grasInMet{Enoncé :}\\
	Soient $A(-1 ;-5) \text { et } \mathrm{B}(2 ; 3)$. Déterminer les coordonnées du vecteur $\vv{A B}$\\
	\grasInMet{Solution :}\\
	Les coordonnées de $\vv{A B}$ sont données par $: x_{B}-x_{A}=2-(-1)=3$ et $y_{B}-y_{A}=3-(-5)=8$. Donc $\Coord{AB}{3}{8}$.\\
	\emph{Pour le vérifier, faire une figure, et vérifier qu'on se déplace bien horizontalement de 3 carreaux vers la droite puis verticalement de 8 carreaux vers le haut afin d'aller du point A au point $B$.}
\end{bmethodet}

%\vspace*{-0.4cm}
\section{Colinéarité de vecteurs dans un repère}

\subsection{Critère de colinéarité}

\begin{bprop}
	Soit $\Coord{u}{x_1}{y_1}$ et $\Coord{v}{x_2}{y_2}$ 2 vecteurs. Alors ces 2 vecteurs sont colinéaires si et seulement si les coordonnées des deux vecteurs sont proportionnelles, soit $x_{1} y_{2}-x_{2} y_{1}=0$
\end{bprop}

\begin{bmethodet}{Montrer que 2 vecteurs sont colinéaires}
	\grasInMet{Enoncé :}
	\begin{enumMet}
		\item Est-ce que les vecteurs $\Coord{u}{-5}{3}$ et $\Coord{v}{15}{-9}$ sont-ils colinéaires ? Justifier.
		\item Est-ce que les vecteurs $\Coord{a}{7}{-3}$  et $\Coord{b}{4}{-2}$ sont-ils colinéaires ? Justifier.
	\end{enumMet}
	\grasInMet{Solution :}\\
	\begin{enumMet}
		\item Puisque $(-5) \times(-9)-3 \times 15=45-45=0$, les vecteurs $\vv{u}$ et $\vv{v}$ sont colinéaires. 
		\item Puisque $7 \times(-2)-(-3) \times 4=-14+12=-2$, les vecteurs $\vv{a}$ et $\vv{b}$ ne sont pas colinéaires.
	\end{enumMet}
\end{bmethodet}

\subsection{Déterminant de deux vecteurs}

\begin{bdefi}
	Soit $\vv{u}$ et $\vv{v}$ deux vecteurs de coordonnées $\left(\begin{array}{l}x \\ y\end{array}\right)$ et $\left(\begin{array}{l}x' \\ y'\end{array}\right)$ dans un repère \Rij.\\
	Le nombre $x y'-y x'$ est appelé \grasInDef{déterminant} des vecteurs $\vv{u}$ et $\vv{v}$.\\
	On note $: \operatorname{det}(\vec{u} ; \vec{v})=\left|\begin{array}{ll}x & x' \\ y & y'\end{array}\right|=x y'-y x'$.
\end{bdefi}

\begin{bprop}
	Dire que $\vv{u}$ et $\vv{v}$ sont colinéaires revient à dire que $\operatorname{det}(\vec{u} ; \vec{v})=0$.
\end{bprop}

\begin{bmethodet}{Vérifier si deux vecteurs sont colinéaires à l'aide du déterminant}
	\grasInMet{Enoncé :}
	Dans chaque cas, vérifier si les vecteurs $\vv{u}$ et $\vv{v}$ sont colinéaires.
	\begin{enumMet}
		\item $\Coord{u}{-6}{10}$ et $\Coord{v}{9}{-15}$
		\item $\Coord{u}{4}{9}$ et $\Coord{v}{11}{23}$
	\end{enumMet}
	\grasInMet{Solution :}\\
	\begin{enumMet}
		\item $\operatorname{det}(\vv{u} ; \vv{v})=\left|\begin{array}{cc}-6 & 9 \\ 10 & -15\end{array}\right|=(-6) \times(-15)-10 \times 9=90-90=0$\\
		Les vecteurs $\vv{u}$ et $\vv{v}$ sont donc colinéaires.
		\item $\operatorname{det}(\vv{u} ; \vv{v})=\left|\begin{array}{ll}4 & 11 \\ 9 & 23\end{array}\right|=4 \times 23-9 \times 11=92-99=-7 \neq 0$\\
		Les vecteurs $\vv{u}$ et $\vv{v}$ ne sont donc pas colinéaires.
	\end{enumMet}
\end{bmethodet}

%\vspace*{-0.4cm}
%\newpage
\subsection{Application de la colinéarité}

\begin{bmethodet}{Appliquer la colinéarité}
	\grasInMet{Enoncé :}
	On considère \Rij un repère du plan. Soit $A(\ -1\ ;\ 1\ ),\ B(\ 3\ ;\ 2\ ),\ C(\ -2\ ;\ -3\ ),\ D(\ 6\ ;\ -1\ ),\ E(\ 5\ ;\ 0\ )$, 5 points du plan.
	\begin{enumMet}
		\item Démontrer que les droites (AB) et (CD) sont parallèles.
		\item Démontrer que les points $\mathrm{E}, \mathrm{B}$ et $\mathrm{D}$ sont alignés.
	\end{enumMet}
	\grasInMet{Solution :}\\
	\begin{enumMet}
		\item $\vv{A B}=\left(\begin{array}{c}3-(-1) \\ 2-1\end{array}\right)=\left(\begin{array}{l}4 \\ 1\end{array}\right)$ et $\vv{C D}=\left(\begin{array}{c}6-(-2) \\ -1-(-3)\end{array}\right)=\left(\begin{array}{l}8 \\ 2\end{array}\right)$.\\
		On remarque que $\vv{C D}=2 \vv{A B}$ donc les coordonnées de $\vv{A B}$ et $\vv{C D}$ sont proportionnelles, et on en déduit que les vecteurs $\vv{A B}$ et $\vv{C D}$ sont colinéaires.\\
		Les droites (AB) et (CD) sont donc parallèles.
		\item $\vv{E B}=\left(\begin{array}{c}3-5 \\ 2-0\end{array}\right)=\left(\begin{array}{c}-2 \\ 2\end{array}\right)$ et $\vv{E D}=\left(\begin{array}{c}6-5 \\ -1-0\end{array}\right)=\left(\begin{array}{c}1 \\ -1\end{array}\right)$\\
		$\operatorname{det}(\vv{E B} ; \vv{E D})=\left|\begin{array}{cc}-2 & 1 \\ 2 & -1\end{array}\right|=-2 \times(-1)-2 \times 1=0$\\
		Les coordonnées de $\vv{E B}$ et $\vv{E D}$ vérifient le critère de colinéarité des vecteurs.
		On en déduit que les vecteurs $\vv{E B}$ et $\vv{E D}$ sont colinéaires.\\
		Les points E, B et D sont donc alignés.
	\end{enumMet}
\end{bmethodet}
\end{document}

